\documentclass[10pt,a4paper]{article}
\usepackage{amsmath,amssymb,amsthm,graphicx}
\usepackage{xcolor}
\usepackage{indentfirst}
\usepackage{amsfonts}
\usepackage{amsmath}
\usepackage{amsthm}
\usepackage{color}
\usepackage[backref=page]{hyperref}
\usepackage[capitalize]{cleveref}
\usepackage{verbatim}
\usepackage[makeroom]{cancel}
\usepackage{rotating}
\usepackage{float}
\usepackage{fullpage}
\usepackage{amssymb}
\usepackage{amsthm}
\usepackage{stmaryrd}
\usepackage{mathrsfs}
\usepackage{epsfig}
\usepackage{indentfirst}
\usepackage{framed}
\usepackage[numbers]{natbib}
\usepackage{color}
\usepackage{graphicx}
\usepackage{float}
\usepackage{bbm}
\usepackage{mathtools}
\usepackage{enumitem}
\usepackage{thmtools}
\usepackage{thm-restate}
\usepackage{longtable}
\usepackage{booktabs}
% \newtheorem{theorem}{Theorem}
\newtheorem{theorem}{Theorem}[section]
\newtheorem{question}[theorem]{Question}
\newtheorem{problem}[theorem]{Problem}
\newtheorem{lemma}[theorem]{Lemma}
\newtheorem{corollary}[theorem]{Corollary}
\newtheorem{definition}[theorem]{Definition}
\newtheorem{proposition}[theorem]{Proposition}
\newtheorem{remark}[theorem]{Remark}
\newtheorem{note}[theorem]{Note}
\newtheorem{fact}[theorem]{Fact}
\newtheorem{claim}[theorem]{Claim}
\newtheorem{assumption}[theorem]{Assumption}
\newtheorem{hypothesis}[theorem]{Hypothesis}
\newtheorem{conjecture}[theorem]{Conjecture}

\crefname{hypothesis}{hypothesis}{hypotheses}

% \numberwithin{theorem}{section}

\newtheorem*{claim*}{Claim}

% Change the parameter below to 0 to remove colors.
\def\colorful{1}

\ifnum\colorful=1
\newcommand{\violet}[1]{{\color{violet}{#1}}}
\newcommand{\orange}[1]{{\color{orange}{#1}}}
\newcommand{\blue}[1]{{\color{blue} {#1}}}
\newcommand{\red}[1]{{\color{red} {#1}}}
\newcommand{\green}[1]{{\color{green} {#1}}}
\newcommand{\pink}[1]{{\color{pink}{#1}}}
\newcommand{\gray}[1]{{\color{gray}{#1}}}

\fi
\ifnum\colorful=0
\newcommand{\violet}[1]{{{#1}}}
\newcommand{\orange}[1]{{{#1}}}
\newcommand{\blue}[1]{{{#1}}}
\newcommand{\red}[1]{{{#1}}}
\newcommand{\green}[1]{{{#1}}}
\fi


\newcommand{\bvee}{\bar{\vee}}
\newcommand{\bwedge}{\bar{\wedge}}
\newcommand{\lsem}{[\![}
\newcommand{\rsem}{]\!]}
% \newcommand{\eqdef}{\stackrel{\rm def}{=}}
\newcommand{\eqdef}{:=}
\newcommand{\defx}[1]{\emph{#1}}

\definecolor{brickred}{rgb}{0.8, 0.25, 0.33}
\definecolor{electriccrimson}{rgb}{1.0, 0.0, 0.25}
\definecolor{harvardcrimson}{rgb}{0.79, 0.0, 0.09}

\definecolor{camouflagegreen}{rgb}{0.47, 0.53, 0.42}
\definecolor{armygreen}{rgb}{0.29, 0.33, 0.13}
\definecolor{darkspringgreen}{rgb}{0.09, 0.45, 0.27}
\definecolor{forestgreen}{rgb}{0.13, 0.55, 0.13}

\definecolor{purpleheart}{rgb}{0.41, 0.21, 0.61}
\newcommand{\inote}[1]{\noindent  \textcolor{blue}{(\textbf{Igor:} #1\noindent)}}
% \newcommand{\bnote}[1]{\noindent \textcolor{purpleheart}{(\textbf{Bruno:} #1\noindent)}}
%\newcommand{\bnote}[1]{}
\newcommand{\bnote}[1]{\noindent  \textcolor{red}{(\textbf{Bruno:} #1\noindent)}}

\newcommand{\XORf}{\mathsf{XOR}^{\mathsf{fnc}}}
\newcommand{\XORg}{\mathsf{XOR}^{\mathsf{graph}}}
\newcommand{\AND}{\mathsf{AND}}
\newcommand{\IPn}{\mathsf{IP}_n}
\newcommand{\binary}{\mathsf{binary}}

% MATH UTILS

\newcommand{\eps}{\varepsilon}
\DeclarePairedDelimiter\ceil{\lceil}{\rceil}
\DeclarePairedDelimiter\floor{\lfloor}{\rfloor}
\newcommand{\one}{\mathbbm{1}}
\DeclareMathOperator{\supp}{supp}
% \DeclareMathOperator{\poly}{poly}
\newcommand{\poly}{\mathsf{poly}}
\DeclareMathOperator{\polylog}{polylog}
\newcommand\inner[2]{\left\langle #1, #2 \right\rangle}
\newcommand{\ovl}{\overline}


% COMPLEXITY CLASSES

% \def\P{{\calp}}
% \def\NP{{\caln\calp}}
% \def\P{{\text{P}}}
% \def\NP{{\text{NP}}}
\newcommand{\complexityfont}[1]{\mathsf{#1}}
\newcommand{\monclass}[1]{\mathsf{m}#1}
% \DeclareMathOperator{\CC}{\mathsf{CC}}
\def\mCC          {\monclass{\CC}}
\def\CC           {\complexityfont{CC}}
\def\TC           {\complexityfont{TC}}
\def\RNC          {\complexityfont{RNC}}
\def\mNC          {\monclass{\NC}}
\def\mTC          {\monclass{\TC}}
\def\EXPSPACE     {\complexityfont{EXPSPACE}}
\def\EXPTIME      {\complexityfont{EXPTIME}
\def\PSPACE       {\complexityfont{PSPACE}}}
\newcommand\NL           {\complexityfont{NL}}
\def\mNL          {\monclass{\NL}}
\def\BPP          {\complexityfont{BPP}}
\def\ZPP          {\complexityfont{ZPP}}
\def\PP           {\complexityfont{PP}}
\def\RP           {\complexityfont{RP}}
\def\coRP         {\complexityfont{coRP}}
\def\P            {\complexityfont{P}}
\def\Ppoly        {\complexityfont{P/poly}}
\def\mPpoly       {\monclass{\Ppoly}}
\def\NP           {\complexityfont{NP}}
\def\coNP         {\complexityfont{coNP}}
\def\EXP          {\complexityfont{EXP}}
\def\AC           {\complexityfont{AC}}
% \def\pL           {\oplus\text{-}\L}
% \def\pL           {\oplus\L}
% \def\pL           {\oplus\complexityfont{L}}
% \def\pL           {\oplus\complexityfont{LOGSPACE}}
\def\pL           {\oplus\complexityfont{L}}
\def\pLpoly       {\pL\complexityfont{/poly}}
\def\NLpoly       {\NL\complexityfont{/poly}}
\def\ACC          {\complexityfont{ACC}}
\def\MAJ          {\complexityfont{MAJ}}
\def\mAC          {\complexityfont{mAC}}
\def\mP           {\monclass{\P}}
\def\NC           {\complexityfont{NC}}
\def\NOT          {\complexityfont{NOT}}
\def\OR           {\complexityfont{OR}}
\def\NAND         {\complexityfont{NAND}}
\def\AND          {\complexityfont{AND}}
\def\SIZE         {\complexityfont{SIZE}}
\def\DEPTH        {\complexityfont{DEPTH}}
\def\mSIZE        {\monclass{\SIZE}}
\def\mDEPTH       {\monclass{\DEPTH}}
\renewcommand{\L} {\complexityfont{L}}
% \renewcommand{\L} {\complexityfont{DLOGSPACE}}
\newcommand{\KW}{\mathop{\text{KW}}}
\newcommand{\mKW}{\mathop{\text{KW}^+}}
\newcommand{\Lclass} {{\normalfont{\textbf{L}}}}
\newcommand{\SPF}[1]{\complexityfont{Span}_{\bbf_{#1}}}
\newcommand{\mSPF}[1]{\monclass{\complexityfont{Span}_{\bbf_{#1}}}}
\newcommand{\semMon}[1]{#1 \cap \monfns}
% \newcommand{\monfns}{\mathsf{mono}}
\newcommand{\monfns}{\mathsf{Mono}}
\newcommand{\GraphPpts}{\mathsf{GraphProperties}}
\newcommand{\HomPres}{\mathsf{HomPreserving}}
\newcommand{\txorsat}{\mathsf{3\mbox{-}XOR\mbox{-}SAT}}
\newcommand{\thornsat}{\mathsf{3\mbox{-}Horn\mbox{-}SAT}}
\newcommand{\tahornsat}{\mathsf{3\mbox{-}AntiHorn\mbox{-}SAT}}
% \newcommand{\txorsat}{\mathsf{3\mbox{-}\!XOR\mbox{-}\!SAT}}
\newcommand{\Gclean}[1]{{#1}_{\mathsf{clean}}}
\newcommand{\spoly}{\mathsf{/poly}}
\newcommand{\hornt}{\calh^3}
\newcommand{\ahornt}{\cala^3}
\newcommand{\NTISP}{\complexityfont{NTiSp}}



% CIRCUIT UTILS

\newcommand{\mDNF}{\monclass{\DNF}}
\newcommand{\mCNF}{\monclass{\CNF}}
\newcommand{\DNF}{\mathsf{DNF}}
\newcommand{\CNF}{\mathsf{CNF}}
\newcommand{\Cmonpref}{\mathsf{C}^{\mathsf{mon}}}
% \newcommand{\Cmon}[1]{\Cmonpref(#1)}
%\newcommand{\Cmon}[1]{\mathsf{mon{\text{-}}size}({#1})}
\newcommand{\Cmon}[1]{\mathsf{mSIZE}({#1})}
\newcommand{\Cman}[1]{\Cmonpref_{\land}(#1)}
\newcommand{\Cskw}[1]{\mathsf{C}_{\mathsf{skew}}(#1)}
\newcommand{\Ccsk}[1]{\mathsf{C}_{\mathsf{skew}}^\circ(#1)}
\newcommand{\Cska}[1]{\mathsf{C}^{\mathsf{skew}}_{\land}(#1)}
\newcommand{\Cszpref}{\mathsf{C}}
% \newcommand{\Cszpref}{\mathsf{C}_\mathsf{size}}
\newcommand{\Csz}[1]{\Cszpref(#1)}
\newcommand{\Csza}[1]{\Cszpref_{\land}(#1)}
% \newcommand{\Csza}[1]{\Cszpref^{\land}(#1)}
\newcommand{\apf}[1]{{#1}^{\cala}}
\newcommand{\GenFml}[1]{\calg_{#1}}
\newcommand{\mdepth}[1]{\mathsf{mDEPTH}({#1})}
\newcommand{\msize}[1]{\Cmon{#1}}

% DISCRETE COMPLEXITY

\newcommand{\bst}{\Gamma}

\newcommand{\gphgn}{\mathcal{G}_N}
\newcommand{\ghclq}{\mathcal{K}_{N,M}}
\newcommand{\boolg}{\mathcal{B}^{\mathsf{mon}}_{n}}
\newcommand{\boolgt}{\mathcal{B}_{2n}}
\newcommand{\mongen}[1]{\mathcal{B}^{\mathsf{mon}}_{#1}}
\newcommand{\blgen}[1]{\mathcal{B}_{#1}}
\newcommand{\Dand}[2]{D_{\cap}({#1} \mid {#2})}
\newcommand{\Dc}[2]{D(#1 \mid #2)}
\newcommand{\Dcg}[1]{\Dc{#1}{\gphgn}}
\newcommand{\Dcb}[1]{\Dc{#1}{\boolgt}}
\newcommand{\Dacg}[1]{\Dand{#1}{\gphgn}}
\newcommand{\Dacb}[1]{\Dand{#1}{\boolgt}}
\newcommand{\row}[1]{R^{#1}}
\newcommand{\col}[1]{C^{#1}}
\newcommand{\rwgc}[1]{\row{#1}_{\ovbr{G}}}
\newcommand{\clgc}[1]{\col{#1}_{\ovbr{G}}}
\newcommand{\flt}[1]{\mathfrak{F}_{\mathsf{mon}}(#1)}
\newcommand{\prj}[1]{[\lvert #1 \rvert]}
\newcommand{\rhobp}{\rho_{\text{BP}}}
\newcommand{\ftdef}[1]{\mathfrak{F}_{\mathsf{def}}(#1)}
\newcommand{\RS}{\mathsf{RS}}

%% COVER GRAPH

% \newcommand{\Vpairs}{\mathsf{Pairs}_U}
% \newcommand{\Vfilters}{\mathsf{Filters}_{A, U, \calb}}
% \newcommand{\Phedges}{\mathcal{E}_{A, U, \calb}}
\newcommand{\Vpairs}{V_\mathsf{pairs}}
\newcommand{\Vfilters}{V_\mathsf{filters}}
\newcommand{\Phedges}{\mathcal{E}}

%%% ABBREVIATIONS
% \renewcommand{\mb}{\bm}
\newcommand{\tu}{\textup}
% \newcommand{\ds}{\displaystyle}
%\newcommand{\ts}{\textstyle}
%\newcommand{\ul}{\underline}
\newcommand{\mr}{\mathrm}
\newcommand{\mi}{\mathit}
\newcommand{\mf}{\mathfrak}
\newcommand{\mc}{\mathcal}
\newcommand{\mb}{\mathbf}
\newcommand{\mbb}{\mathbb}
\newcommand{\mds}{\mathds}
% \newcommand{\scr}{\mathscr}
%\newenvironment{pf}{\begin{proof}}{\end{proof}}

% MATH OPERATORS

\def\sseq{\subseteq}
\newcommand{\set}[1]{\left\{#1\right\}}
\newcommand{\overbar}[1]{\mkern 1.5mu\overline{\mkern-1.5mu#1\mkern-1.5mu}\mkern 1.5mu}
\newcommand{\ovbr}{\overbar}
\def\sm{\setminus}
\def\var{\mathop{\text{\rm var}}\nolimits}
\newcommand{\ds}{\oplus}
\newcommand{\card}[1]{\lvert#1\rvert}
\newcommand{\abs}[1]{\left\lvert#1\right\rvert}
\newcommand{\hwt}[1]{\card{#1}_1}
% \newcommand{\cmpl}[1]{\bar{#1}}
% \newcommand{\cmpl}[1]{\ovbr{#1}}
% \newcommand{\cmpl}[1]{{#1}^{\mathsf{c}}}
\newcommand{\cmpl}[1]{#1^{\mathsf{c}}}
\newcommand{\brkt}[1]{\left\langle #1 \right\rangle}

% MATH OPERATORS (MATRIX)

% \newcommand{\tT}{\intercal}
% \newcommand{\tT}{\top}
% \newcommand{\tT}{T}
\newcommand{\tT}{\mathsf{T}}
\newcommand{\tr}[1]{#1^\tT}
\newcommand{\Diagld}[1]{\Diag \left( \ldm{A} \right)}

% BOOLEAN

%\newcommand{\CNF}{\mathsf{CNF}}
\newcommand{\blt}{\set{0,1}}
\newcommand{\vecp}{\vec{p}}
\newcommand{\vecq}{\vec{q}}
\newcommand{\vecr}{\vec{r}}
% \newcommand{\power}[1]{2^{#1}}
\newcommand{\power}[1]{\calp({#1})}
\DeclareMathOperator{\DT}{DT}
% \newcommand{\dtdepth}[1]{\DT_\text{\footnotesize depth}(#1)}
\newcommand{\dtdepth}[1]{\mathsf{DT}_\mathsf{depth}(#1)}
% \newcommand{\dtsize}[1]{\DT_\text{\footnotesize\rm size}(#1)}
\newcommand{\dtsize}[1]{\mathsf{DT}_\mathsf{size}(#1)}

% BOOLEAN FUNCTIONS

% \DeclareMathOperator{\clique}{Clique}
\newcommand{\clique}{\mathsf{Clique}}
\DeclareMathOperator{\cycle}{CYCLE}
\DeclareMathOperator{\thr}{THR}
% \newcommand{\vcprm}{k\mbox{-}\!\mathsf{Vertex\mbox{-}Cover}}
\newcommand{\kvc}{k\mbox{-}\!\mathsf{Vertex\mbox{-}Cover}}
\newcommand{\clqprm}[1]{#1\mbox{-}\!\clique}
\newcommand{\kclq}{k\mbox{-}\!\clique}
% \newcommand{\kclq}{\mathsf{CLIQUE}(n,k)}
% \newcommand{\kclq}{\mathsf{Clique}(n,k)}
\newcommand{\kmtc}{k\mbox{-}\!\Matching}
\newcommand{\Matching}{\mathsf{Matching}}
\newcommand{\BPM}{\mathsf{Bipartite}\text{-}\mathsf{Matching}}
\newcommand{\Parity}{\mathsf{Parity}}
\newcommand{\maj}{\mathsf{Majority}}
\newcommand{\MOD}{\mathsf{MOD}}
\newcommand{\andp}{\mathsf{POLY(q,s)}}
% \newcommand{\kclq}{\mathsf{Clique}_{k,n}}
\newcommand{\kcycl}{k$-$\cycle}
\newcommand{\MCSP}{\mathsf{MCSP}}
\newcommand{\stconn}{\mathsf{ST\text{-}CONN}}
\newcommand{\ustconn}{\mathsf{U\text{-}ST\text{-}CONN}}
\newcommand{\conn}{\mathsf{Connectivity}}
\newcommand{\GEN}{\mathsf{GEN}}
\newcommand{\SAT}{\mathsf{SAT}}
\newcommand{\dSAT}{2\text{-}\mathsf{SAT}}
\newcommand{\MtxSs}[1]{\llbracket#1\rrbracket}
\newcommand{\MSP}{\mathsf{MSP}}
\newcommand{\xor}{\oplus}
\newcommand{\bigxor}{\bigoplus}
\newcommand{\OddFactor}{\mathsf{OddFactor}}
\newcommand{\EqOne}{\mathsf{EqualONE}}
\newcommand{\EvenCC}{\mathsf{EvenCC}}
\newcommand{\BipEvenCC}{\mathsf{Bipartite}\text{-}\mathsf{EvenCC}}
\newcommand{\BipOddFactor}{\mathsf{Bipartite}\text{-}\OddFactor}


% CSP

\newcommand{\CSPSAT}{\mathsf{CSP}\text{-}\mathsf{SAT}}
% \newcommand{\COQf}[1]{\mathsf{COQ}\text{-}\mathsf{SAT}_{#1}}
\newcommand{\COQf}[1]{\mathsf{COQ}_{#1}}
\newcommand{\CSPf}[1]{\mathsf{CSP}\text{-}\mathsf{SAT}_{#1}}
% \newcommand{\CSPfn}[2]{\mathsf{CSP}\text{-}\mathsf{SAT}_{#1}^{#2}}
\newcommand{\CSPfn}[2]{\CSPf{#1}^{#2}}
\newcommand{\COQfn}[3]{\mathsf{COQ}_{#1}^{(#2,#3)}}
\newcommand{\CSP}{\mathsf{CSP}}
\newcommand{\COQ}{\mathsf{COQ}}
\newcommand{\Pol}{\mathsf{Pol}}
\newcommand{\Inv}{\mathsf{Inv}}
\newcommand{\ccln}[1]{\brkt{#1}}
\newcommand{\Cstr}[1]{\mathsf{Constraints}(#1)}
\newcommand{\Cstrn}[2]{\mathsf{Constraints}_{#1}(#2)}
\newcommand{\proj}{\mathsf{proj}}
\newcommand{\dual}{\mathsf{dual}}
\newcommand{\acmred}{\leq^{\AC^0}_m\!}
\newcommand{\acmcmpl}{$\acmred$\!-complete}
\newcommand{\mnlmred}{\leq^{\mNL}_m}
\newcommand{\mprojred}{\leq^{\mathsf{mProj}}_m}
\newcommand{\mormred}{\leq^{\mathsf{mOR}}_m}
\newcommand{\clonefont}[1]{\mathrm{#1}}


% COMMUNICATION COMPLEXITY

\newcommand{\PLS}{\mathsf{PLS}}
\newcommand{\Root}{\emptyset}

% LIFTING

% \newcommand{\SrchF}[1]{\mathsf{Search}_{#1}}
\newcommand{\SrchF}[1]{\mathsf{Search}(#1)}
\newcommand{\Clauses}[1]{\mathsf{Clauses}(#1)}


% CIRCUIT CLASSES

% \newcommand{\AC}{\mathsf{AC}}
% \newcommand{\NC}{\mathsf{NC}}

% FUSION METHOD

\DeclareMathOperator{\Gen}{\mathsf{GEN}}

% PROBABILITY STUFF

% \newcommand{\probabilityfont}[1]{\mathsf{#1}}
\newcommand{\Bin}{\text{Bin}}
\newcommand{\dstr}{\sim}
\newcommand{\probabilityfont}[1]{\mb{#1}}
\newcommand{\rnd}[1]{\probabilityfont{#1}}
\def\Pr{\mathop{{}\mathbb{P}}}
% \def\Pr{\mathop{\text{\rm Pr}}}
% \def\Exp{\mathop{\bbe}limits}
% \def\Exp{\mathop{\bbe}}
\def\Exp{\mathop{{}\mathbb{E}}}
\def\var{\mathop{\text{\rm var}}\nolimits}
\def\pcritk{p_{\rm c}^{K_k}}
\def\pcritl{p_{\rm c}^{C_k}}
\newcommand{\inuf}{\in_U}
\newcommand{\cktsz}{\calc}
% \newcommand{\inuf}{\in_calu}
\newcommand{\rndY}{\probabilityfont{Y}}
\newcommand{\rndN}{\probabilityfont{N}}
\newcommand{\rndW}{\probabilityfont{W}}
\newcommand{\rndP}{\probabilityfont{P}}
\newcommand{\rndS}{\probabilityfont{S}}
\newcommand{\rndT}{\probabilityfont{T}}
\newcommand{\rndX}{\probabilityfont{X}}
\newcommand{\rndG}{\probabilityfont{G}}
\newcommand{\rndH}{\probabilityfont{H}}
\newcommand{\rndF}{\probabilityfont{F}}
\newcommand{\rndK}{\probabilityfont{K}}
\newcommand{\rndc}{\probabilityfont{c}}
\newcommand{\rndx}{\probabilityfont{x}}
\newcommand{\flws}{\sim}


\def\bba{{\mathbb A}}
\def\bbb{{\mathbb B}}
\def\bbc{{\mathbb C}}
\def\bbd{{\mathbb D}}
\def\bbe{{\mathbb E}}
\def\bbf{{\mathbb F}}
\def\bbg{{\mathbb G}}
\def\bbh{{\mathbb H}}
\def\bbi{{\mathbb I}}
\def\bbj{{\mathbb J}}
\def\bbk{{\mathbb K}}
\def\bbl{{\mathbb L}}
\def\bbm{{\mathbb M}}
\def\bbn{{\mathbb N}}
\def\bbo{{\mathbb O}}
\def\bbp{{\mathbb P}}
\def\bbq{{\mathbb Q}}
\def\bbr{{\mathbb R}}
\def\bbs{{\mathbb S}}
\def\bbt{{\mathbb T}}
\def\bbu{{\mathbb U}}
\def\bbv{{\mathbb V}}
\def\bbw{{\mathbb W}}
\def\bbx{{\mathbb X}}
\def\bby{{\mathbb Y}}
\def\bbz{{\mathbb Z}}

\def\cala{{\mathcal A}}
\def\calb{{\mathcal B}}
\def\calc{{\mathcal C}}
\def\cald{{\mathcal D}}
\def\cale{{\mathcal E}}
\def\calf{{\mathcal F}}
\def\calg{{\mathcal G}}
\def\calh{{\mathcal H}}
\def\cali{{\mathcal I}}
\def\calj{{\mathcal J}}
\def\calk{{\mathcal K}}
\def\call{{\mathcal L}}
\def\calm{{\mathcal M}}
\def\caln{{\mathcal N}}
\def\calo{{\mathcal O}}
\def\calp{{\mathcal P}}
\def\calq{{\mathcal Q}}
\def\calr{{\mathcal R}}
\def\cals{{\mathcal S}}
\def\calt{{\mathcal T}}
\def\calu{{\mathcal U}}
\def\calv{{\mathcal V}}
\def\calw{{\mathcal W}}
\def\calx{{\mathcal X}}
\def\caly{{\mathcal Y}}
\def\calz{{\mathcal Z}}

\def\scra{{\mathscr A}}
\def\scrb{{\mathscr B}}
\def\scrc{{\mathscr C}}
\def\scrd{{\mathscr D}}
\def\scre{{\mathscr E}}
\def\scrf{{\mathscr F}}
\def\scrg{{\mathscr G}}
\def\scrh{{\mathscr H}}
\def\scri{{\mathscr I}}
\def\scrj{{\mathscr J}}
\def\scrk{{\mathscr K}}
\def\scrl{{\mathscr L}}
\def\scrm{{\mathscr M}}
\def\scrn{{\mathscr N}}
\def\scro{{\mathscr O}}
\def\scrp{{\mathscr P}}
\def\scrq{{\mathscr Q}}
\def\scrr{{\mathscr R}}
\def\scrs{{\mathscr S}}
\def\scrt{{\mathscr T}}
\def\scru{{\mathscr U}}
\def\scrv{{\mathscr V}}
\def\scrw{{\mathscr W}}
\def\scrx{{\mathscr X}}
\def\scry{{\mathscr Y}}
\def\scrz{{\mathscr Z}}
\newcommand{\headsubsec}[1]{\paragraph{#1.}}
\newcommand{\headsubsecline}[1]{\headsubsec{#1}}
\newcommand{\gaphyp}{monotone gap conjecture}
\newcommand{\gaphyps}{monotone gap conjectures}
\newcommand{\ckthyp}{monotone circuit gap conjecture}
\newcommand{\fmlhyp}{monotone formula gap conjecture}
\title{Non-monotone circuit lower bounds from efficient monotone simulations}
\author{Bruno P. Cavalar, Igor C. Oliveira}
\date{doi:10.4230/LIPIcs.CCC.2023.29; CC BY4.0}
\begin{document}
\maketitle
\setcounter{section}{2}
\noindent\textit{Editorial note.} Original authors: Bruno P. Cavalar and Igor C. Oliveira. arXiv:2305.06821v1 (2023-05-11), https://arxiv.org/abs/2305.06821v1. Creative Commons Attribution4.0.\par\medskip
\noindent\textit{Editorial note.} License: https://creativecommons.org/licenses/by/4.0/\par\medskip
\noindent\textit{Editorial note.} Selected complete native author theorem/proof/context and full bibliography in standalone layout. Native input/mathhead and mathematical style definitions flattened as explicit spans; unused frontmatter and figure packages omitted. No mathematical changes or new proofs. Original CC BY4.0 license retained.\par\medskip
\noindent\textit{Allowed standard prerequisites (editorial).}\begin{itemize}
\item Known monotone lower bounds, logspace algorithms and threshold circuit bounds. Exact author prior theorem statements retained with citations. Full author depth-reduction/padding support proofs included.
\end{itemize}\medskip
\section{Preliminaries}

\subsection{Notation}

\subparagraph{Boolean functions.}

We denote by $\monfns$ the set of all monotone Boolean functions.
We define $\poly = \set{n \mapsto n^C : C \in \bbn}$.
A Boolean function
    $f : \blt^{\binom{n}{2}} \to \blt$ is said to be a graph property
    if $f(G) = f(H)$ for any two isomorphic graphs $G$ and $H$.
Let $\calf = \set{f_n}_{n \in \bbn}$
    be a sequence of graph properties, where $f_n$ is defined over undirected graphs on
    $n$ vertices.
    We say that $\calf$ is
    \emph{preserved under homomorphisms} if,
    whenever there is a homomorphism from a graph $G$ to a graph $H$,
    we have $\calf(G) \leq \calf(H)$.
    We denote by $\HomPres$ 
    the set of all 
    graph properties which are preserved under homomorphisms.
    Note that $\HomPres \sseq \monfns$.

    \subparagraph{Boolean circuits.}

    We denote by $\AC^0_d[s]$ 
    the family of Boolean functions computed by
    size-$s$,
    depth-$d$
    Boolean circuits with unbounded fan-in $\set{\land,\lor}$-gates and input literals from $\{x_1, \overline{x_1}, \ldots, x_n, \overline{x_n}\}$.
    We write $\AC^0[s]$ as a shorthand for
    $\bigcup_{d = 1}^\infty \AC^0_d[s]$,
    and $\AC^0$ as a shorthand of $\AC^0[n^{O(1)}] = \AC^0[\poly]$.
    We will also refer to $\AC^0_d[\poly]$
    by $\AC^0_d$.
We write $\DNF[s]$ to denote the family of Boolean functions
    computed by size-$s$ DNFs, where size is measured by  number of terms.
    We write $\CNF[s]$ analogously.
    We write $\SIZE[s]$ to denote
    the family of Boolean functions
    computed by size-$s$ circuits.
    We write $\DEPTH[d]$ to denote
    the family of Boolean functions
    computed by fan-in 2 circuits of depth $d$.
We denote by $\AC^0[\xor]$ the family of Boolean functions computed by
polynomial-size $\AC^0$ circuits with unbounded fan-in $\xor$-gates.

We denote by $\L$ the 
family of
Boolean functions
computed by  logspace machines,
and by $\NL$ the family of Boolean functions
computed by polynomial-time nondeterministic logspace machines.
Moreover, we denote by $\pL$
the family of Boolean functions computed by
polynomial-time nondeterministic logspace machines with a \emph{parity}
acceptance condition (i.e.,
an input is accepted if the number of accepting paths is odd).

\subparagraph{Circuit complexity.}

Given a circuit class $\calc$, we write
    $\monclass{C}$ to denote the
    monotone version of $\calc$.
Given a function $f$, we write $\Cmon{f}$
    to denote the size of the smallest monotone circuit computing
    $f$ and $\mdepth{f}$ to denote the smallest depth of a fan-in 2
    monotone circuit computing $f$.
Given two Boolean functions $f,g$, we write
$f \mprojred g$ if
there exists a many-one reduction from $f$ to $g$
in which each bit of the reduction is a monotone
projection\footnote{A monotone projection is a projection without
negations.}
of the input.

\subparagraph{Miscellanea.}

Let $\alpha \in \blt^n$ .
    We 
    define
    $\hwt{\alpha} := \sum_{i=1}^n \alpha_i$.
    We call $\card{\alpha}_1$ the \defx{Hamming weight} of $\alpha$.
    We let
    $\supp(\alpha) = \set{i \in [n] : \alpha_i = 1}$.
We let $\thr_{k,n} : \blt^n \to \blt$ be the Boolean function
    such that
    $\thr_{k,n}(x) = 1 \iff \hwt{x} \geq k$.



\subsection{Background results}


The next lemma, which is proved via a standard ``guess-and-verify''
approach, shows that nondeterministic logspace computations can be
simulated by circuits of size $2^{n^{\varepsilon}}$ and of depth $d =
O_\varepsilon(1)$.

\begin{lemma}[Folklore; see, e.g.,~{\cite[Lemma 8.1]{AHMPS08}}]
    \label{lemma:ac0_nc1_sim}
    For all $\eps > 0$, we have
    $\NL \sseq \AC^0[2^{n^\eps}]$.
\end{lemma}







\begin{theorem}[Polynomial-size constant-depth vs.~larger monotone depth]
    \label{thm:ac0-fml-graph-intro}
    For every $k \geq 1$,
    we have
    $\semMon{\AC^0} \not\sseq \mDEPTH[(\log n)^k]$. Moreover, this separation holds for a monotone graph property.
\end{theorem}

\begin{restatable}[Polynomial-size constant-depth vs.~larger monotone size]{theorem}{acxormp}
    \label{thm:ac02-mp-intro}
    For every $k \geq 1$,
    we have
$\semMon{\AC^0[\oplus]} \not\sseq \mSIZE[2^{(\log n)^k}]$.
\end{restatable}

\begin{theorem}[\cite{gkrs19}]
    \label{thm:xor-sat}
    There exists $\eps > 0$ 
    such that 
    $\semMon{\pL} \not\sseq \mSIZE[2^{o(n^{\eps})}]$.
    Moreover, 
    this separation is witnessed by $\txorsat$.
\end{theorem}

\begin{theorem}
    \label{thm:ac0-fml-graph}
    For every $i \geq 1$,
    we have
    $\semMon{\AC^0} \cap \GraphPpts \not\sseq \mDEPTH[(\log n)^i]$.
    In particular, we have
    $\semMon{\AC^0} \cap \GraphPpts \not\sseq \mNC^i$.
\end{theorem}

\begin{theorem}[New circuit lower bounds from monotone simulations]
    \label{thm:simul_consequences_intro}
    There exists $\eps > 0$
    such that the following holds.
    \begin{enumerate}
        \item 
            If 
            $\semMon{\AC^0_3} 
            \subseteq \mNC^1$,
            then 
            $\NP 
            \not\sseq
            \AC^0_3[2^{o(n)}]$.
        \item 
            If 
            $\semMon{\AC^0_4} 
            \subseteq \mSIZE[\poly]$,
            then 
            $\NP 
            \not\sseq
            \AC^0_4[ 2^{o( \sqrt{n}/\log n )} ]$.
        \item If 
            $\semMon{\AC^0}  
            \sseq 
            \mSIZE[\poly]$,
            then
            $\NC^2 
            \not\sseq
            \NC^1$.
        \item 
            If $\semMon{\NC^1} \sseq \mSIZE[2^{O(n^\eps)}]$,
            then 
            $\NC^2 
            \not\sseq
            \NC^1$.
        \item 
            If 
            $\semMon{\AC^0} \cap \GraphPpts \subseteq \mSIZE[\mathsf{poly}]$,
            then $\NP \not\sseq \NC^1$.
        \item 
            If $\semMon{\NC^1} \cap \GraphPpts \sseq \mSIZE[\mathsf{poly}]$,
            then 
            $\L
            \not\sseq
            \NC^1$.
    \end{enumerate}
\end{theorem}

\begin{lemma}[Folklore; see, e.g.,~\cite{AKRRV01, DBLP:conf/coco/OliveiraS019}]
    \label{lemma:pl-to-ac02}
    Let $f : \blt^n \to \blt$
    be a Boolean function computed by
    a $\pL$ machine.
    For every $\delta > 0$,
    there exists an $\AC^0[\oplus]$ circuit
    of size $2^{n^\delta}$ that computes $f$.
\end{lemma}
\begin{proof}
    Let $M$ be a $\pL$-machine computing $f$.
    Without loss of generality, we may assume that
    each configuration in the 
    configuration graph $G$ of $M$
    is \emph{time-stamped} -- in other words,
    each configuration 
    carries the information of the number of
    computational steps it takes to arrive at it.\footnote{Formally,
        we can define a $\pL$-machine $M'$ such that
        the configurations of $M'$
        are
        $(C,t)$, where
        $C$ is a configuration of $M$, and $t = 0,1,\dots,m = n^{O(1)}$
        is a number denoting the time in which the configuration was
        achieved.
        A configuration $(C,t)$ can only reach a configuration
        $(C',t+1)$ in the configuration graph of $M'$.
    }
    We may also assume that
    every accepting computation takes exactly the same amount of time,
    which means that every path from the starting configuration
    $v_{\mathsf{start}}$
    to the accepting configuration
    $v_{\mathsf{accept}}$
    has the same length in the configuration graph.
    These assumptions imply that the configuration graph is \emph{layered} 
    (because a configuration with time-stamp $t$
        can only 
        point to
    configurations with time-stamp $t+1$)
    and acyclic.
    Note that, for a fixed machine,
    the configuration graph can be computed from the input string using a projection.  

    Let $m = n^{O(1)}$ be the time that an accepting computation takes.
    We now show
    how to count (modulo 2)
    the number of accepting paths from
    $v_{\mathsf{start}}$
    to
    $v_{\mathsf{accept}}$ with a depth-$d$ $\AC^0[\oplus]$ circuit.
    First, choose
    $m^{1/d}-1$ configurations $v_1,\dots,v_{m^{1/d}-1}$
    (henceforth called ``checkpoints'')
    from $V(G)$, such that the configuration
    $v_i$ is at the level
    $i \cdot m^{1-1/d}$ in the configuration graph
    (i.e., it takes $i \cdot m^{1-1/d}$ time steps to arrive at $v_i$).
    For convenience, we let $v_0 = v_{\mathsf{start}}$
    and $v_{m^{1/d}} = v_{\mathsf{accept}}$.
    We then count the number of paths from 
    from $v_{\mathsf{start}}$
    to $v_{\mathsf{accept}}$
    that go through $v_1,\dots,v_{m^{1/d}-1}$,
    and sum over all possible choices of the checkpoints.
    Since the graph is layered and
    each path from $v_0$ to $v_{m^{1/d}}$ has length exactly $m$,
    there is only one choice of checkpoints that witnesses
    a given path from $v_0$ to $v_{m^{1/d}}$,
    so no path is counted twice in this summation.
    Letting
    $\#\mathsf{paths}(s,t,\ell)$ denote the number of paths between
    configurations $s$ and $t$ with distance exactly $\ell$,
    we obtain
    \begin{equation*}
        \#\mathsf{paths}(v_{0},v_{m^{1/d}}, m)
        =
        \sum_{v_1,\dots,v_{m^{1/d}-1}}
        \prod_{i=0}^{m^{1/d}-1}
        \#\mathsf{paths}(v_i,v_{i+1}, m^{1-1/d}).
    \end{equation*}
    The above calculation can be done in modulo 2 with an unbounded fan-in XOR gate
    (replacing the summation)
    and an unbounded fan-in AND gate
    (replacing the product).
    Note that the formula above is recursive.
    Repeating the same computation for calculating (modulo 2)
    the expression
    $\#\mathsf{paths}(v_i,v_{i+1}, m^{1-1/d})$ for each $i$,
    we obtain a depth-$2d$ $\AC^0[\oplus]$ circuit
    for calculating the number of paths from $v_{\mathsf{start}}$ to
    $v_{\mathsf{accept}}$ (modulo 2).
    Clearly, the total size of the circuit is
    $2^{O(m^{1/d} \cdot \log m)}$,
    which is smaller than $2^{n^\delta}$ for a large enough constant $d$.
\end{proof}

First, we recall a result
of~\cite{DBLP:journals/combinatorica/BabaiGW99},
which proves monotone lower bounds for the following function.
Let 
$\OddFactor_n : \blt^{\binom{n}{2}} \to \blt$
be the function that accepts a given 
graph
if 
it
contains an \emph{odd factor} -- in other words,
a spanning subgraph in which the degree of every vertex is odd.
Babai, Gál and Wigderson~\cite{DBLP:journals/combinatorica/BabaiGW99}
proved
the following result:
\begin{theorem}[\cite{DBLP:journals/combinatorica/BabaiGW99}]
    \label{thm:bip-oddfactor-lb}
        Any monotone formula computing 
        $\OddFactor_n$ 
        has size $2^{\Omega(n)}$,
        and any monotone circuit computing
        $\OddFactor_n$ 
        has size $n^{\Omega(\log n)}$.
\end{theorem}

The proof in \cite{DBLP:journals/combinatorica/BabaiGW99}  is actually for the case of
\emph{bipartite graphs}, but it easily extends to general graphs, since
the bipartite case reduces to the general case by a monotone projection.
The formula lower bound stated above is slightly stronger because it makes
use of asymptotically optimal lower bounds on the randomized communication
complexity of $\mathsf{DISJ}_n$
\citep{DBLP:journals/siamdm/KalyanasundaramS92}, which were not available
to \cite{DBLP:journals/combinatorica/BabaiGW99}. We remark that, with a
different language,
a monotone circuit lower bound for $\OddFactor$ is also implicitly proved
in Feder and
Vardi~\cite[Theorem 30]{DBLP:journals/siamcomp/FederV98}.



We now recall an upper bound for $\OddFactor$,  implicitly proved
in an unpublished note due to Johannsen~\cite{joh03}.

\begin{theorem}[\cite{joh03}]
    \label{thm:oddfactor-logspace}
    We have $\OddFactor \in \L$.
\end{theorem}
\begin{proof}
    We first recall the following observation about the $\OddFactor$
    function, which appears in different forms in the literature
    (see~\cite[Lemma 4.1]{DBLP:journals/jacm/Urquhart87}
    or~\cite[Lemma 18.16]{jukna_2012};
    see also~\cite[Proposition 1]{joh03} for a different proof.)
    \begin{claim*}
        A graph $G$ has an odd factor if and only if
        every connected component of $G$ has an even number of vertices.
    \end{claim*}
    \begin{proof}
        If a graph $G$ has an odd factor, 
        we can conclude
        that every
        connected
        component of $G$ has an even number of vertices
        from the well-known observation
        that in every graph 
        there is an even number of vertices of odd degree.

        Now suppose that every connected component of $G$ 
        has an even number of vertices.
        We will iteratively construct an odd factor $F$ of $G$.
        We begin with the empty graph.
        We take any two vertices $u,v$ in the same
        connected component of $G$
        which currently have even degree in $F$,
        and consider any path
        $P = (x_1,\dots,x_{k})$
        between $u$ and $v$,
        where $x_1 = u$ and $x_{k}=v$.
        If the edge $x_i x_{i+1}$ is currently in $F$, we remove 
        $x_i x_{i+1}$ from $F$; otherwise, we add $x_i x_{i+1}$ to
        $F$.
        It's easy to check that, in every iteration of this procedure,
        only the vertices $u$ and $v$ have the parity of their degree
        changed in $F$; the degree of every other vertex stays the same
        (modulo 2).
        Since every connected component has an even number of vertices,
        this means that, eventually, every vertex in $F$ will have odd
        degree.
    \end{proof}
    Now it's easy to check in logspace if every connected component of $G$ has an even
    number of vertices using Reingold's algorithm for undirected
    connectivity~\cite{DBLP:conf/stoc/Reingold05}.
    It suffices to check if, for every vertex $v$ of $G$,
    the number of vertices reachable from $v$ is odd.
\end{proof}

Now, if we only desire to obtain a function in $\AC^0$
not computed by monotone circuits of depth $(\log n)^i$,
we can follow the same argument of~\Cref{thm:ac02-mp-intro},
using \Cref{lemma:ac0_nc1_sim} instead of \Cref{lemma:pl-to-ac02}.
In order to obtain moreover a monotone \emph{graph property} witnessing
this separation, we will need the following lemma,
which enables us to obtain a graph property after ``padding'' a graph property.
We defer the proof of this lemma to
the end of this section.

\begin{restatable}{lemma}{paddinggraph}
    \label{lemma:padding_graph}
    Let $f : \blt^{\binom{n}{2}} \to \blt$
    be a monotone graph property on graphs of $n$ vertices.
    The following holds.
    \begin{enumerate}
        \item If $f \in \NC^i$ for some $i > 1$, then
            there exists a monotone graph property 
            $g$ on graphs of 
            $N = 2^{(\log n)^i}$ vertices
            such that 
            $g \in \NC^1$
            and 
            $f \mprojred g$.
        \item If $f \in \NL$, 
            then
            for all $\eps > 0$
            there exists
            a monotone graph property 
            $g$ on graphs of 
            $N = 2^{n^\eps}$ vertices
            such that 
            $g$ can be computed by  $\AC^0$ circuits of size $N^{2+o(1)}$
            and 
            $f \mprojred g$.
        \item If $f \in \pL$, 
            then
            for all $\eps > 0$
            there exists
            a monotone graph property 
            $g$ on graphs of 
            $N = 2^{n^\eps}$ vertices
            such that 
            $g$ can be computed by  $\AC^0[\xor]$ circuits of size $N^{2+o(1)}$
            and 
            $f \mprojred g$.
    \end{enumerate}
\end{restatable}

We are now ready to prove
\Cref{thm:ac0-fml-graph}.
\begin{proof}[Proof of \Cref{thm:ac0-fml-graph}]
            Fix $n \in \bbn$ and take an $\eps < 1/i$.
            Observing that $\L \sseq \NL$,
            from
            \Cref{thm:oddfactor-logspace}
            and item (2) of \Cref{lemma:padding_graph}
            we conclude that
            there exists a monotone graph property $f$
            on 
            $N = 2^{n^\eps}$
            vertices
            such that 
            $f \in \AC^0$
            and $\OddFactor_n \mprojred f$.
            By \Cref{thm:bip-oddfactor-lb},
            any monotone circuit computing $f$
            has depth 
            $\Omega(n) = \Omega((\log N)^{1/\eps}) \gg (\log N)^i$.
\end{proof}

Raz and Wigderson~\cite{RW92} observed that there exists a monotone function
$f \in \NC^1 \sm \mNC$.
Using~\Cref{lemma:padding_graph},
we observe moreover that it's possible to obtain this separation
with a monotone graph property.

\begin{proposition}
    \label{prop:nc1-mnc-graph}
    We have $\semMon{\NC^1} \cap \GraphPpts \not\sseq \mNC$.
\end{proposition}
\begin{proof}
    Observing that $\L \sseq \NC^2$,
    we conclude 
    from
    \Cref{thm:oddfactor-logspace}
    and item (1) of \Cref{lemma:padding_graph}
    that
    there exists a monotone graph property $f$
    on 
    $N = 2^{(\log n)^2}$ 
    vertices
    such that 
    $f \in \NC^1$
    and 
    $\OddFactor_n \mprojred f$.
    By 
    \Cref{thm:bip-oddfactor-lb},
    any monotone circuit computing $f$
    has depth 
    $\Omega(n) = \Omega(2^{\sqrt{\log N}})$,
    which implies $f \not\in \mNC$.
    \qedhere
\end{proof}


\subsection{Efficient monotone padding for graph properties}

We will now prove \Cref{lemma:padding_graph}.
We first recall some low-depth circuits
for computing threshold functions, which we will use to design a circuit for efficiently computing
the adjacency matrix of induced subgraphs.





\begin{theorem}[\cite{hwwy94}]
    \label{thm:polylog-thr-ac0}
Let $d > 0$ be a constant.
The function
$\thr_{(\log n)^d,n}$ 
can be computed by an $\AC^0$ circuit of size $n^{o(1)}$
and depth $d+O(1)$.
\end{theorem}

\begin{theorem}[\cite{aks83}]
    \label{thm:thr-nc1}
    For every $k \in [n]$,
    the function $\thr_{k,n}$ can be computed by a
    circuit of depth $O(\log n)$ and size
    $n^{O(1)}$.
\end{theorem}


\begin{lemma}
    \label{lemma:induced_subgraph_circuit}
    There exists a circuit $C^k_n$ with $\binom{n}{2} + n$ inputs
    and $\binom{k}{2}$ outputs which,
    when given as input
    an adjacency matrix 
    of a graph $G$ on $n$ vertices
    and a characteristic vector of a set $S \sseq [n]$
    such that $\card{S} \leq k$,
    outputs the adjacency matrix of the graph $G[S]$,
    padded with isolated vertices when $\card{S} < k$.
    The circuit has constant-depth and size $n^{2+o(1)}$
    when $k = \polylog(n)$, and
    size $n^{O(1)}$ and depth $O(\log n)$ otherwise.
\end{lemma}
\begin{proof}
    Let $\set{x_{ij}}_{i,j \in [n]}$ encode the adjacency matrix of $G$.
    Let $\alpha \in \blt^n$ be the characteristic vector of $S$.
    Let $i,j \in [k]$.
    Note that $\set{i,j} \in E(G[S])$
    if and only if
    there exists
    $a,b \in [n]$
    such that
    \begin{itemize}
        \item 
            $\alpha_a$ is the $i$-th non-zero entry of $\alpha$,
        \item 
            $\alpha_b$ is the $j$-th non-zero entry of $\alpha$, and
        \item 
            $x_{ab} = 1$ (i.e., $a$ and $b$ are connected in $G$).
    \end{itemize}
    We first consider the case $k = \polylog(n)$.
    In this case,
    the first two conditions can be checked with 
    circuits of size
    $n^{o(1)}$
    using \Cref{thm:polylog-thr-ac0}.
    Therefore, we can 
    compute 
    if $i$ and $j$ are adjacent
    using $n^{2+o(1)}$ gates and constant depth.
    As there are at most 
    $(\log n)^{O(1)}$ 
    such pairs,
    we can output $G[S]$
    with at most 
    $n^{2+o(1)}$ gates.

    For any $k$, the first two conditions can be checked with an $\NC^1$
    circuit by Theorem~\ref{thm:thr-nc1}.
    Since there are at most $n^2$ pairs $i,j$, the entire adjacency matrix
    can be computed with a $O(\log n)$-depth and polynomial-size circuit.
\end{proof}

We are ready to prove \Cref{lemma:padding_graph}.


\begin{proof}[Proof of \Cref{lemma:padding_graph}]
    We first prove (1). 
    Fix $n \in \bbn$
    and let $N = 2^{(\log n)^i}$.
    For a graph $G$ on $N$
    vertices such that 
    $\card{E(G)} \leq \binom{n}{2}$,
    let
    $\Gclean{G}$ be the graph 
    obtained
    from $G$ by removing isolated vertices from $G$ one-by-one,
    in lexicographic order,
    until one of the following two 
    conditions are satisfied:
    (1) there are no more isolated vertices in $\Gclean{G}$,
    \emph{or}
    (2) $\Gclean{G}$ has exactly $n$ vertices.
    Let $g : \blt^{\binom{N}{2}} \to \blt$
    be the monotone graph property defined as follows:
    \begin{align*}
        g(G) := 
        \left(
            \card{E(G)} > 
            \binom{n}{2}
        \right)
        \lor
        \left(
            \card{V(\Gclean{G})} > n
        \right)
        \lor
            (f(\Gclean{G}) = 1).
    \end{align*}
    Note that $g$ accepts a graph $G$
    if and only if at least 
    one of the following three conditions are satisfied:
    \begin{enumerate}
        \item 
            $G$ has at most $\binom{n}{2}$ edges,
            $\Gclean{G}$ has exactly $n$ vertices
            and 
            $f(\Gclean{G}) = 1$,
            or
        \item 
            $G$ has more than $\binom{n}{2}$ edges,
            or
        \item 
            $\Gclean{G}$ has more than $n$ vertices.
    \end{enumerate}
    We observe that the monotonicity of $g$ follows from 
    the monotonicity of $f$.
    We also claim that $g$ is a graph property.
    Indeed, the graph $\Gclean{G}$ is the same (up to isomorphism),
    irrespective  of the order according to which the isolated vertices are
    removed from $G$. Moreover, the function $f$ is also a graph property.
    Because of this, all the three conditions above
    are preserved under isomorphisms.

    We first observe that $f$ is a monotone projection of $g$.
    Indeed, given a graph $G$ on $n$ vertices,
    we can easily construct by a monotone
    projection a graph $G'$ on $N$ vertices
    and at most $\binom{n}{2}$ edges
    such that
    $f(G) = g(G')$.
    We just let $G'$ have a planted copy of $G$,
    and all other vertices are isolated.
    Then $\Gclean{G'} = G$ (up to isomorphism)
    and $g(G') = f(\Gclean{G}) = f(G)$.

    We now show how to compute $g$ in $\NC^1$.
    Let $\set{x_{ij}}_{i,j \in [N]}$ be the input bits of $g$,
    corresponding to the adjacency matrix of a graph $G$.
    The circuit computes as follows.

    \begin{enumerate}
        \item If $\card{E(G)} > \binom{n}{2}$, accept the graph $G$.
        \item Compute the characteristic vector $\alpha \in \blt^N$
            of the set of all non-isolated vertices of $G$.
            If $\hwt{\alpha} > n$, accept the graph $G$.
        \item Compute $\Gclean{G}$ and output $f(\Gclean{G})$.
    \end{enumerate}


    Note that checking if
    $\card{E(G)} > \binom{n}{2}$
    can be done 
    in $\NC^1$ by Theorem~\ref{thm:thr-nc1}.
    Moreover,
    for all $i \in [N]$,
    we have $\alpha_i = \bigvee_{j \in [N]} x_{ij}$,
    and therefore $\alpha_i$ can be computed by a circuit of depth $O(\log N)$ 
    and $O(N)$ gates.
    In total, 
    the vector $\alpha$ can be computed
     with $O(N^2)$ gates and $O(\log N)$ depth.
    Finally, we can check if
    $\hwt{\alpha} > n$
    in $\NC^1$ with a threshold circuit.

    For the final step, we compute $\Gclean{G}$.
    If $\hwt{\alpha} = n$,
    note that $\Gclean{G} = G[\supp(\alpha)]$.
    When $\hwt{\alpha} < n$,
    then $\Gclean{G}$ is $G[\supp(\alpha)]$ padded with isolated vertices.
    We can therefore compute $\Gclean{G}$
    with the circuit 
    $C_N^n$
    of Lemma~\ref{lemma:induced_subgraph_circuit}.
    Moreover, since $f \in \NC^i$,
    we have that $f$ can be computed by a circuit of size
    $n^{O(1)} = N^{o(1)}$
    and depth
    $O((\log n)^i) = O(\log N)$.
    Therefore, computing $f(\Gclean{G})$ can be done in $\NC^1$. Overall, we get that $g \in \NC^1$.

    In order to prove (2), it suffices to modify the proof above.
    The modification can be briefly described as follows.
    We let $N = 2^{n^\eps}$.
    Every time \Cref{lemma:induced_subgraph_circuit} is applied,
    we use the $\AC^0$ circuit instead of the $\NC^1$ circuit,
    since $n = \polylog(N)$. This ammounts to $N^{2+o(1)}$ many gates with
    unbounded fan-in.
    Moreover, since by assumption 
    $f \in \NL$,
    applying \Cref{lemma:ac0_nc1_sim} we obtain 
    an
    $\AC^0$ circuit for
    $f$ of size
    $2^{n^{\eps/2}} = N^{o(1)}$,
    so we can compute $f(\Gclean{G})$ in constant depth
    with $N^{o(1)}$ gates.

    Finally, for (3) it suffices to apply the same argument used for (2),
    replacing an application of~\Cref{lemma:ac0_nc1_sim}
    by an application of \Cref{lemma:pl-to-ac02}.
\end{proof}

\subsection{Non-monotone lower bounds
from monotone simulations}


We now show that if monotone circuits are able to
efficiently simulate non-monotone circuits computing monotone Boolean
functions,
then striking complexity separations follow.
We also show a result of this kind for simulations of graph properties.
We first prove a lemma connecting the simulation of $\AC^0$ circuits
with the simulation of $\NL$ machines.

\begin{lemma}
    \label{thm:nc1-ac0-size}
    For all constants $\eps > 0$ and $C \geq 1$, 
    if 
    $\semMon{\AC^0} 
    \sseq \mSIZE[2^{O((\log n)^C)}]$, 
    then
    $\semMon{\NL} \sseq \mSIZE[2^{o(n^\eps)}]$.
\end{lemma}
\begin{proof}
    We prove the contrapositive.
    Suppose that there exists $\eps > 0$
    such that
    $\semMon{\NL} \not\sseq \mSIZE[2^{o(n^\eps)}]$.
    This means that there exists a monotone function $f$
    such that 
    $f \in \NL$
    and any monotone circuit computing $f$ has size $2^{\Omega(n^\eps)}$.

    Let $\delta = \eps/(2C)$ and let $m = 2^{n^{\delta}}$.
    Let $g : \blt^{n} \times \blt^{m} \to \blt$
    be the Boolean function defined as
    $g(x,y) = f(x)$.
    Note that $g$ is a function on $N := m + n = 2^{\Theta(n^{\delta})}$
    bits.
    By \Cref{lemma:ac0_nc1_sim},
    there exists an $\AC^0$ circuit computing
    $f$ of size 
    $2^{n^\delta} = N^{O(1)}$.
    Moreover, any monotone circuit computing $g$
    has size
    $2^{\Omega(n^\eps)} = 2^{\Omega((\log N)^{\eps / \delta})} =
    2^{\Omega((\log N)^{2C})}$.
\end{proof}




Next, we recall the strongest known monotone circuit 
and formula
lower 
bounds
for a
monotone function in $\NP$. 

\begin{theorem}[\cite{DBLP:conf/stoc/PitassiR17}]
    \label{thm:pr17-fnp}
     $\semMon{\NP} \not\sseq \mDEPTH[o(n)]$.
\end{theorem}

\begin{theorem}[\cite{CKR20}]
    \label{thm:ckr20}
     
    $\semMon{\NP} \not\sseq \mSIZE[2^{o(\sqrt{n}/\log n)}]$.
\end{theorem}


We are now ready to state and prove our first result regarding
new complexity separations from monotone simulations. Recall that obtaining explicit lower bounds against depth-$3$ $\AC^0$  circuits of size $2^{\omega(n^{1/2})}$ is a major challenge in circuit complexity theory, while the best lower bound on the size of depth-4 $\AC^0$ circuits
computing a function in $\NP$ is
currently $2^{\Omega(n^{1/3})}$~\cite{DBLP:conf/stoc/Hastad86}.
Moreover, no strict separation is known in the following sequence of inclusions
of complexity classes:
$\ACC \sseq \TC^0 \sseq \NC^1 \sseq \L \sseq \NL \sseq \pL \sseq \NC^2$.
We show that 
efficient monotone simulations would bring new results in both of these
fronts. 
(We stress that all lower bound consequences appearing below refer to
separations against non-uniform circuits.)\footnote{In other words, all upper bounds are
\emph{uniform}, but the lower bounds hold even for \emph{non-uniform}
circuits. Note that this is stronger than lower bounds for uniform
circuits.}
 

\begin{theorem}
    \label{cor:xor-sat-ag}
    Let $\calc$ be 
    a class of circuits.
    There exists $\eps > 0$ such that
    the following holds:
    \begin{enumerate}
        \item 
            \label{item:ac03}
            If $\semMon{\AC^0_3} \subseteq \mNC^1$,
            then 
            $\NP 
            \not\sseq
            \AC^0_3[2^{o(n)}]$.
        \item 
            \label{item:ac04}
            If 
            $\semMon{\AC^0_4} \subseteq \mSIZE[\poly]$,
            then 
            $\NP 
            \not\sseq
            \AC^0_4[ 2^{o( \sqrt{n}/\log n )} ]$.
        \item 
            \label{item:subexp}
            If $\semMon{\calc} \sseq \mSIZE[2^{O(n^\eps)}]$,
            then 
            $\NC^2 
            \not\sseq
            \calc$.
        \item 
            \label{item:ac-poly}
            If 
            $\semMon{\AC^0} 
            \sseq 
            \mSIZE[\poly]$,
            then
            $\NC^2 
            \not\sseq
            \NC^1$.
    \end{enumerate}
\end{theorem}
\begin{proof}
    We will prove each item separately.
    \begin{enumerate}[label=\emph{Proof of (\arabic*)}., leftmargin=*]
        \item 
            Let us assume that
            $\semMon{\AC^0_3} \subseteq \mNC^1$.
            Let $f$ be the function of Theorem~\ref{thm:pr17-fnp}.
            For a contradiction,  suppose that 
            $f \in 
            \AC^0_3[2^{o(n)}]$.
            Let $\alpha : \bbn \to \bbn$ be such that
            $\alpha(n) \to_n \infty$
            and
            $f$ has a depth-3 $\AC^0$ circuit
            of size
            $2^{n/\alpha}$.
            Let $m = 2^{n/(10\cdot\alpha)}$
            and
            let $g : \blt^n \times \blt^m \to \blt$
            be the function
            $g(x,y) = f(x)$.
            Let $N = n+m = (1+o(1))2^{n/(10\cdot\alpha)}$.
            Clearly, the function $g$ has 
            a depth-3 $\AC^0$ circuit 
            of size $2^{n/\alpha} = N^{O(1)}$.
            Since $g$ is monotone,
            we conclude from the assumption 
            that $g$ is computed by a polynomial-size monotone formula.
            Now, since $f(x) = g(x,1^{m})$,
            we obtain a monotone formula of size
            $N^{O(1)} = 2^{o(n)}$ for computing $f$,
             which contradicts the lower bound of
             Theorem~\ref{thm:pr17-fnp}.
        \item 
            Similar to the proof of item~(\ref{item:ac03}),
            but using \Cref{thm:ckr20} instead.
        \item 
            Suppose that $\NC^2 \sseq \calc$.
            By \Cref{thm:xor-sat}, there exists a monotone function $f \in \NC^2$
            on $n$ bits
            and a number $\eps > 0$
            such that
            $f \notin \mSIZE[2^{o(n^{\eps})}]$.
            Therefore, for any $\delta > 0$ such that $\delta < \eps$,
            we have $f \notin \mSIZE[2^{O(n^\delta)}]$.
            Since, by assumption, we have $f \in \NC^2 \sseq \calc$,
            we obtain
            $\semMon{\calc} \not\sseq \mSIZE[2^{O(n^\delta)}]$.
        \item 
            If $\NC^2 \sseq \NC^1$, then,
            by 
            item (\ref{item:subexp}),
            we get
            $\semMon{\NC^1} \not\sseq \mSIZE[2^{o(n^\eps)}]$.
            From \Cref{thm:nc1-ac0-size},
            we obtain
            $\semMon{\AC^0} 
            \not\sseq \mSIZE[\poly]$.
            \qedhere
    \end{enumerate}
\end{proof}



As a motivation to 
the ensuing discussion,
we recall a result of
Rossman~\cite{rossman_2008},
who 
showed that any homomorphism-preserving
graph property computed by \(\AC^0\) circuits is also computed
by monotone \(\AC^0\) circuits.

\begin{theorem}[\cite{rossman_2008}]
    \label{t:ros08}
    $\AC^0
    \cap 
    \HomPres 
    \sseq \mDNF[\poly]$.
\end{theorem}

This inspires the question of whether general graph properties
can also be efficiently simulated by monotone circuits.
We show that, if true,
such simulations would imply strong complexity separations.
Let us first recall an exponential monotone circuit lower bound for
monotone graph properties, and we will be ready to state and prove our main
result.


\begin{theorem}[\cite{DBLP:journals/combinatorica/AlonB87}]
    \label{thm:ab-graph}
    There exists $\eps > 0$
    such that
    $\semMon{\NP} \cap \GraphPpts \not\sseq \mSIZE[2^{o(n^\eps)}]$.
\end{theorem}





\begin{theorem}
    \label{cor:matching-ag}
    Let $\calc$ be 
    a class of circuits.
    The following holds:
    \begin{enumerate}
        \item 
            \label{item:gpsize}
            If $\semMon{\calc} \cap \GraphPpts \sseq \mSIZE[\poly]$,
            then 
            $\L
            \not\sseq
            \calc$.
        \item 
            \label{item:gpdepth}
            If 
            $\semMon{\calc} \cap \GraphPpts \sseq 
            \mDEPTH[o(\sqrt{n})]$, where $n$ denotes the number of input bits,
            then
            $\L
            \not\sseq
            \calc$.
        \item 
            \label{item:kclq}
            If 
            $\semMon{\AC^0} \cap \GraphPpts \subseteq \mSIZE[\poly]$,
            then $\NP \not\sseq \NC^1$.
    \end{enumerate}
\end{theorem}
\begin{proof}
    We will prove each item separately.

    \begin{enumerate}[label=\emph{Proof of (\arabic*)}., leftmargin=*]
        \item 
            Suppose that $\L \sseq \calc$.
            By \Cref{thm:bip-oddfactor-lb},
            the monotone graph property 
            $\OddFactor$
            satisfies
            $\OddFactor \notin \mSIZE[\poly]$.
            Moreover, we have the upper bound $\OddFactor \in \L$ by
            \Cref{thm:oddfactor-logspace}.
            Since, by assumption, we have $\OddFactor \in \L \sseq \calc$,
            we obtain
            $\semMon{\calc} \cap \GraphPpts \not\sseq \mSIZE[\poly]$.
        \item 
            Suppose that $\L \sseq \calc$.
            By 
            \Cref{thm:oddfactor-logspace,thm:bip-oddfactor-lb}, 
            there exists a monotone graph property $f \in \L$ such
            that $f \notin \mDEPTH[o(\sqrt{n})]$.
            Since, by assumption, we have
            $f \in \L \sseq \calc$,
            we obtain
            $\semMon{\calc} \cap \GraphPpts \not\sseq \mDEPTH[o(\sqrt{n})]$.

        \item 
            Suppose that $\NP \sseq \NC^1$.
            By \Cref{thm:ab-graph}, there exists a monotone graph property
            $f \in \NC^1$ such that
            $\Cmon{f} = 2^{\Omega(n^\eps)}$ for some $\eps > 0$.
            Let $\delta = \eps/2$.
            By Lemma~\ref{lemma:padding_graph} (Item $2$),
            there exists a monotone graph property $g$ on
            $N = 2^{n^\delta}$ vertices
            computed by
            an $\AC^0$ circuit of size $N^{2+o(1)}$
            such that $f$ is a monotone projection of $g$.
            \Cref{thm:ab-graph}
            implies that 
            any monotone circuit computing $f$ has size
            $2^{\Omega(n^{\eps})} 
            = 2^{\Omega((\log N)^2)} = N^{\omega(1)}$.
            \qedhere
    \end{enumerate}
\end{proof}

\newcommand{\etalchar}[1]{$^{#1}$}
\begin{thebibliography}{HWWY94}

\bibitem[AB87]{DBLP:journals/combinatorica/AlonB87}
Noga Alon and Ravi~B. Boppana.
\newblock The monotone circuit complexity of boolean functions.
\newblock {\em Combinatorica}, 7(1):1--22, 1987.

\bibitem[ABI{\etalchar{+}}09]{DBLP:journals/jcss/AllenderBISV09}
Eric Allender, Michael Bauland, Neil Immerman, Henning Schnoor, and Heribert
  Vollmer.
\newblock The complexity of satisfiability problems: Refining {S}chaefer's
  theorem.
\newblock {\em J. Comput. Syst. Sci.}, 75(4):245--254, 2009.

\bibitem[AG87]{ajtai_gurevich_1987}
Mikl\'{o}s Ajtai and Yuri Gurevich.
\newblock Monotone versus positive.
\newblock {\em J. Assoc. Comput. Mach.}, 34(4):1004--1015, 1987.

\bibitem[AHM{\etalchar{+}}08]{AHMPS08}
Eric Allender, Lisa Hellerstein, Paul McCabe, Toniann Pitassi, and Michael~E.
  Saks.
\newblock Minimizing disjunctive normal form formulas and {AC}$^0$ circuits
  given a truth table.
\newblock {\em {SIAM} J. Comput.}, 38(1):63--84, 2008.

\bibitem[AK11]{akiyama_factors_2011}
Jin Akiyama and Mikio Kano.
\newblock {\em Factors and {Factorizations} of {Graphs}}, volume 2031 of {\em
  Lecture {Notes} in {Mathematics}}.
\newblock Springer Berlin Heidelberg, Berlin, Heidelberg, 2011.

\bibitem[AKR{\etalchar{+}}01]{AKRRV01}
Eric Allender, Michal Kouck{\'{y}}, Detlef Ronneburger, Sambuddha Roy, and
  V.~Vinay.
\newblock Time-space tradeoffs in the counting hierarchy.
\newblock In {\em Proceedings of the 16th Annual {IEEE} Conference on
  Computational Complexity, Chicago, Illinois, USA, June 18-21, 2001}, pages
  295--302. {IEEE} Computer Society, 2001.

\bibitem[AKS83]{aks83}
Mikl{\'{o}}s Ajtai, J{\'{a}}nos Koml{\'{o}}s, and Endre Szemer{\'{e}}di.
\newblock An {$O(n \log n)$} sorting network.
\newblock In {\em Proceedings of the 15th Annual {ACM} Symposium on Theory of
  Computing, 25-27 April, 1983, Boston, Massachusetts, {USA}}, pages 1--9.
  {ACM}, 1983.

\bibitem[And85]{andreev1985method}
Alexander~E Andreev.
\newblock On a method for obtaining lower bounds for the complexity of
  individual monotone functions.
\newblock {\em Doklady Akademii Nauk SSSR}, 282:1033--1037, 1985.

\bibitem[BCRV03]{playing_one}
Elmar Böhler, Nadia Creignou, Steffen Reith, and Heribert Vollmer.
\newblock Playing with boolean blocks, part i: Post’s lattice with
  applications to complexity theory.
\newblock {\em SIGACT News}, 2003.

\bibitem[BCRV04]{playing_two}
Elmar Böhler, Nadia Creignou, Steffen Reith, and Heribert Vollmer.
\newblock Playing with boolean blocks, part ii: Constraint satisfaction
  problems.
\newblock {\em ACM SIGACT-Newsletter}, 35, 2004.

\bibitem[BDHM92]{DBLP:journals/mst/BuntrockDHM92}
Gerhard Buntrock, Carsten Damm, Ulrich Hertrampf, and Christoph Meinel.
\newblock Structure and importance of logspace-mod class.
\newblock {\em Math. Syst. Theory}, 25(3):223--237, 1992.

\bibitem[Ber82]{berkowitz}
S.~J. Berkowitz.
\newblock On some relationships between monotone and non-monotone circuit
  complexity.
\newblock Technical Report. University of Toronto, 1982.

\bibitem[BGW99]{DBLP:journals/combinatorica/BabaiGW99}
L{\'{a}}szl{\'{o}} Babai, Anna G{\'{a}}l, and Avi Wigderson.
\newblock Superpolynomial lower bounds for monotone span programs.
\newblock {\em Combinatorica}, 19(3):301--319, 1999.

\bibitem[BHST14]{DBLP:conf/icalp/BlaisHST14}
Eric Blais, Johan H{\aa}stad, Rocco~A. Servedio, and Li{-}Yang Tan.
\newblock On {DNF} approximators for monotone boolean functions.
\newblock In {\em International Colloquium on Automata, Languages, and
  Programming \emph{(ICALP)}}, pages 235--246, 2014.

\bibitem[BL88]{DBLP:conf/crypto/Leichter88}
Josh~Cohen Benaloh and Jerry Leichter.
\newblock Generalized secret sharing and monotone functions.
\newblock In {\em Advances in Cryptology \emph{(CRYPTO)}}, pages 27--35, 1988.

\bibitem[BST13]{BST13}
Eric Blais, Dominik Scheder, and Li-Yang Tan.
\newblock Ajtai-gurevich redux. {M}anuscript.
\newblock 2013.

\bibitem[BT96]{DBLP:journals/jacm/BshoutyT96}
Nader~H. Bshouty and Christino Tamon.
\newblock On the {F}ourier spectrum of monotone functions.
\newblock {\em J. {ACM}}, 43(4):747--770, 1996.

\bibitem[Bul17]{DBLP:conf/focs/Bulatov17}
Andrei~A. Bulatov.
\newblock A dichotomy theorem for nonuniform csps.
\newblock In Chris Umans, editor, {\em 58th {IEEE} Annual Symposium on
  Foundations of Computer Science, {FOCS} 2017, Berkeley, CA, USA, October
  15-17, 2017}, pages 319--330. {IEEE} Computer Society, 2017.

\bibitem[Bul18]{DBLP:journals/siglog/Bulatov18}
Andrei~A. Bulatov.
\newblock Constraint satisfaction problems: complexity and algorithms.
\newblock {\em {ACM} {SIGLOG} News}, 5(4):4--24, 2018.

\bibitem[CDM21]{cdm21}
Arkadev Chattopadhyay, Rajit Datta, and Partha Mukhopadhyay.
\newblock Lower bounds for monotone arithmetic circuits via communication
  complexity.
\newblock In {\em {STOC} '21: 53rd Annual {ACM} {SIGACT} Symposium on Theory of
  Computing, Virtual Event, Italy, June 21-25, 2021}, pages 786--799. {ACM},
  2021.

\bibitem[CHO{\etalchar{+}}20]{choprs20}
Lijie Chen, Shuichi Hirahara, Igor~C. Oliveira, J{\'a}n Pich, Ninad Rajgopal,
  and Rahul Santhanam.
\newblock {Beyond Natural Proofs: Hardness Magnification and Locality}.
\newblock In {\em 11th Innovations in Theoretical Computer Science Conference
  (ITCS 2020)}, volume 151, pages 70:1--70:48, 2020.

\bibitem[CKR20]{CKR20}
Bruno~Pasqualotto Cavalar, Mrinal Kumar, and Benjamin Rossman.
\newblock Monotone circuit lower bounds from robust sunflowers.
\newblock In {\em {LATIN} 2020: Theoretical Informatics - 14th Latin American
  Symposium, S{\~{a}}o Paulo, Brazil, January 5-8, 2021, Proceedings}, volume
  12118 of {\em Lecture Notes in Computer Science}, pages 311--322. Springer,
  2020.

\bibitem[CKS01]{DBLP:books/daglib/0004131}
Nadia Creignou, Sanjeev Khanna, and Madhu Sudan.
\newblock {\em Complexity classifications of Boolean constraint satisfaction
  problems}, volume~7 of {\em {SIAM} monographs on discrete mathematics and
  applications}.
\newblock {SIAM}, 2001.

\bibitem[COS17]{cos15}
Xi~Chen, Igor~C. Oliveira, and Rocco~A. Servedio.
\newblock Addition is exponentially harder than counting for shallow monotone
  circuits.
\newblock In {\em S{TOC}'17---{P}roceedings of the 49th {A}nnual {ACM} {SIGACT}
  {S}ymposium on {T}heory of {C}omputing}, pages 1232--1245. ACM, New York,
  2017.

\bibitem[CZ16]{DBLP:conf/stoc/ChattopadhyayZ16}
Eshan Chattopadhyay and David Zuckerman.
\newblock Explicit two-source extractors and resilient functions.
\newblock In {\em Symposium on Theory of Computing \emph{(STOC)}}, pages
  670--683, 2016.

\bibitem[dRGR22]{DBLP:journals/sigact/RezendeGR22}
Susanna~F. de~Rezende, Mika G{\"{o}}{\"{o}}s, and Robert Robere.
\newblock Guest column: Proofs, circuits, and communication.
\newblock {\em {SIGACT} News}, 53(1):59--82, 2022.

\bibitem[FV98]{DBLP:journals/siamcomp/FederV98}
Tom{\'{a}}s Feder and Moshe~Y. Vardi.
\newblock The computational structure of monotone monadic {SNP} and constraint
  satisfaction: {A} study through datalog and group theory.
\newblock {\em {SIAM} J. Comput.}, 28(1):57--104, 1998.

\bibitem[GGLR98]{DBLP:conf/focs/GoldreichGLR98}
Oded Goldreich, Shafi Goldwasser, Eric Lehman, and Dana Ron.
\newblock Testing monotonicity.
\newblock In {\em Symposium on Foundations of Computer Science, \emph{(FOCS)}},
  pages 426--435, 1998.

\bibitem[GI12]{DBLP:journals/toc/GoldreichI12}
Oded Goldreich and Rani Izsak.
\newblock Monotone circuits: One-way functions versus pseudorandom generators.
\newblock {\em Theory Comput.}, 8(1):231--238, 2012.

\bibitem[GJW18]{gjw18}
Mika G{\"{o}}{\"{o}}s, Rahul Jain, and Thomas Watson.
\newblock Extension complexity of independent set polytopes.
\newblock {\em {SIAM} J. Comput.}, 47(1):241--269, 2018.

\bibitem[GKRS19]{gkrs19}
Mika G\"{o}\"{o}s, Pritish Kamath, Robert Robere, and Dmitry Sokolov.
\newblock Adventures in monotone complexity and {TFNP}.
\newblock In {\em 10th {I}nnovations in {T}heoretical {C}omputer {S}cience},
  volume 124 of {\em LIPIcs. Leibniz Int. Proc. Inform.}, pages Art. No. 38,
  19. Schloss Dagstuhl. Leibniz-Zent. Inform., Wadern, 2019.

\bibitem[GMOR15]{DBLP:conf/tcc/GuoMOR15}
Siyao Guo, Tal Malkin, Igor~C. Oliveira, and Alon Rosen.
\newblock The power of negations in cryptography.
\newblock In {\em Theory of Cryptography Conference \emph{(TCC)}}, pages
  36--65, 2015.

\bibitem[GS92]{gs92}
Michelangelo Grigni and Michael Sipser.
\newblock Monotone complexity.
\newblock In {\em Proceedings of the London Mathematical Society Symposium on
  Boolean Function Complexity}, page 57–75, USA, 1992. Cambridge University
  Press.

\bibitem[H{\aa}s86]{DBLP:conf/stoc/Hastad86}
Johan H{\aa}stad.
\newblock Almost optimal lower bounds for small depth circuits.
\newblock In {\em Proceedings of the 18th Annual {ACM} Symposium on Theory of
  Computing, May 28-30, 1986, Berkeley, California, {USA}}, pages 6--20, 1986.

\bibitem[Hir22]{Hir22}
Shuichi Hirahara.
\newblock {NP}-hardness of learning programs and partial {MCSP}.
\newblock In {\em Symposium on Foundations of Computer Science \emph{(FOCS)}},
  2022.

\bibitem[HWWY94]{hwwy94}
Johan H{\aa}stad, Ingo Wegener, Norbert Wurm, and Sang{-}Zin Yi.
\newblock Optimal depth, very small size circuits for symmetric functions in
  {AC}$^0$.
\newblock {\em Inf. Comput.}, 108(2):200--211, 1994.

\bibitem[IKL{\etalchar{+}}19]{DBLP:journals/jacm/IkenmeyerKLLMS19}
Christian Ikenmeyer, Balagopal Komarath, Christoph Lenzen, Vladimir Lysikov,
  Andrey Mokhov, and Karteek Sreenivasaiah.
\newblock On the complexity of hazard-free circuits.
\newblock {\em J. {ACM}}, 66(4):25:1--25:20, 2019.

\bibitem[JCG97]{10.1145/263867.263489}
Peter Jeavons, David Cohen, and Marc Gyssens.
\newblock Closure properties of constraints.
\newblock {\em J. ACM}, 44(4):527–548, jul 1997.

\bibitem[Jea98]{DBLP:journals/tcs/Jeavons98}
Peter Jeavons.
\newblock On the algebraic structure of combinatorial problems.
\newblock {\em Theor. Comput. Sci.}, 200(1-2):185--204, 1998.

\bibitem[JLL76]{DBLP:journals/mst/JonesLL76}
Neil~D. Jones, Y.~Edmund Lien, and William~T. Laaser.
\newblock New problems complete for nondeterministic log space.
\newblock {\em Math. Syst. Theory}, 10:1--17, 1976.

\bibitem[Joh03]{joh03}
Jan Johannsen.
\newblock The complexity of satisfiability problems with two occurrences.
\newblock Unpublished Manuscript, 2003.

\bibitem[Juk12]{jukna_2012}
Stasys Jukna.
\newblock {\em Boolean Function Complexity - Advances and Frontiers}.
\newblock Springer, 2012.

\bibitem[Kra97]{Krajicek97}
Jan Kraj\'{\i}\v{c}ek.
\newblock Interpolation theorems, lower bounds for proof systems, and
  independence results for bounded arithmetic.
\newblock {\em J. Symbolic Logic}, 62(2):457--486, 1997.

\bibitem[KS92]{DBLP:journals/siamdm/KalyanasundaramS92}
Bala Kalyanasundaram and Georg Schnitger.
\newblock The probabilistic communication complexity of set intersection.
\newblock {\em {SIAM} J. Discret. Math.}, 5(4):545--557, 1992.

\bibitem[Kup21]{DBLP:conf/lics/Kuperberg21}
Denis Kuperberg.
\newblock Positive first-order logic on words.
\newblock In {\em Symposium on Logic in Computer Science \emph{(LICS)}}, pages
  1--13, 2021.

\bibitem[Kup22]{DBLP:journals/corr/abs-2201-11619}
Denis Kuperberg.
\newblock Positive first-order logic on words and graphs.
\newblock {\em CoRR}, abs/2201.11619, 2022.

\bibitem[KW93]{DBLP:conf/coco/KarchmerW93}
Mauricio Karchmer and Avi Wigderson.
\newblock On span programs.
\newblock In {\em Structure in Complexity Theory Conference \emph{(CCC)}},
  pages 102--111, 1993.

\bibitem[LT09]{DBLP:journals/tcs/LaroseT09}
Beno{\^{\i}}t Larose and Pascal Tesson.
\newblock Universal algebra and hardness results for constraint satisfaction
  problems.
\newblock {\em Theor. Comput. Sci.}, 410(18):1629--1647, 2009.

\bibitem[Oko82]{okolnishnikova_1982}
E.~A. Okol'nishnikova.
\newblock The effect of negations on the complexity of realization of monotone
  {B}oolean functions by formulas of bounded depth.
\newblock {\em Metody Diskret. Analiz.}, (38):74--80, 1982.

\bibitem[Oli15]{Oliveira15}
Igor~C. Oliveira.
\newblock Unconditional lower bounds in complexity theory, 2015.
\newblock (PhD Thesis, Columbia University).

\bibitem[OSS19]{DBLP:conf/coco/OliveiraS019}
Igor~C. Oliveira, Rahul Santhanam, and Srikanth Srinivasan.
\newblock Parity helps to compute majority.
\newblock In {\em 34th Computational Complexity Conference, {CCC} 2019, July
  18-20, 2019, New Brunswick, NJ, {USA}}, pages 23:1--23:17, 2019.

\bibitem[Pos41]{post41}
Emil~L. Post.
\newblock {\em The Two-Valued Iterative Systems of Mathematical Logic. (AM-5)}.
\newblock Princeton University Press, 1941.

\bibitem[PR17]{DBLP:conf/stoc/PitassiR17}
Toniann Pitassi and Robert Robere.
\newblock Strongly exponential lower bounds for monotone computation.
\newblock In {\em Symposium on Theory of Computing \emph{(STOC)}}, pages
  1246--1255, 2017.

\bibitem[Pud97]{DBLP:journals/jsyml/Pudlak97}
Pavel Pudl{\'{a}}k.
\newblock Lower bounds for resolution and cutting plane proofs and monotone
  computations.
\newblock {\em J. Symb. Log.}, 62(3):981--998, 1997.

\bibitem[Qui53]{quine_1953}
W.~V. Quine.
\newblock Two theorems about truth functions.
\newblock {\em Bol. Soc. Mat. Mexicana}, 10:64--70, 1953.

\bibitem[Raz85a]{razborov_boolean_85}
A.~A. Razborov.
\newblock Lower bounds on the monotone complexity of some {B}oolean functions.
\newblock {\em Dokl. Akad. Nauk SSSR}, 281(4):798--801, 1985.

\bibitem[Raz85b]{razborov1985lower}
Alexander~A Razborov.
\newblock Lower bounds on monotone complexity of the logical permanent.
\newblock {\em Mathematical Notes of the Academy of Sciences of the USSR},
  37(6):485--493, 1985.

\bibitem[Rei05]{DBLP:conf/stoc/Reingold05}
Omer Reingold.
\newblock Undirected st-connectivity in log-space.
\newblock In Harold~N. Gabow and Ronald Fagin, editors, {\em Proceedings of the
  37th Annual {ACM} Symposium on Theory of Computing, Baltimore, MD, USA, May
  22-24, 2005}, pages 376--385. {ACM}, 2005.

\bibitem[Rez23]{susanna_2023}
Susanna Rezende.
\newblock Personal communication, 2023.

\bibitem[RM99]{DBLP:journals/combinatorica/RazM99}
Ran Raz and Pierre McKenzie.
\newblock Separation of the monotone {NC} hierarchy.
\newblock {\em Combinatorica}, 19(3):403--435, 1999.

\bibitem[Ros08a]{rossman_2008}
Benjamin Rossman.
\newblock Homomorphism preservation theorems.
\newblock {\em J. ACM}, 55(3):Art. 15, 53, 2008.

\bibitem[Ros08b]{DBLP:conf/stoc/Rossman08}
Benjamin Rossman.
\newblock On the constant-depth complexity of $k$-clique.
\newblock In {\em Symposium on Theory of Computing \emph{(STOC)}}, pages
  721--730, 2008.

\bibitem[Ros17a]{rossman_entropy}
Benjamin Rossman.
\newblock An entropy proof of the switching lemma and tight bounds on the
  decision-tree size of {AC}$^0$.
\newblock 2017.

\bibitem[Ros17b]{rossman_16}
Benjamin Rossman.
\newblock An improved homomorphism preservation theorem from lower bounds in
  circuit complexity.
\newblock In {\em Innovations in Theoretical Computer Science Conference
  \emph{(ITCS)}}, pages 27:1--27:17, 2017.

\bibitem[RPRC16]{rprc16}
Robert Robere, Toniann Pitassi, Benjamin Rossman, and Stephen~A. Cook.
\newblock Exponential lower bounds for monotone span programs.
\newblock In {\em {IEEE} 57th Annual Symposium on Foundations of Computer
  Science, {FOCS} 2016, 9-11 October 2016, Hyatt Regency, New Brunswick, New
  Jersey, {USA}}, pages 406--415. {IEEE} Computer Society, 2016.

\bibitem[RW92]{RW92}
Ran Raz and Avi Wigderson.
\newblock Monotone circuits for matching require linear depth.
\newblock {\em J. {ACM}}, 39(3):736--744, 1992.

\bibitem[Sch78]{DBLP:conf/stoc/Schaefer78}
Thomas~J. Schaefer.
\newblock The complexity of satisfiability problems.
\newblock In {\em Proceedings of the 10th Annual {ACM} Symposium on Theory of
  Computing, May 1-3, 1978, San Diego, California, {USA}}, pages 216--226.
  {ACM}, 1978.

\bibitem[Sto95]{DBLP:conf/lics/Stolboushkin95}
Alexei~P. Stolboushkin.
\newblock Finitely monotone properties.
\newblock In {\em Symposium on Logic in Computer Science \emph{(LICS)}}, pages
  324--330, 1995.

\bibitem[Tar88]{tardos_1988}
\'{E}. Tardos.
\newblock The gap between monotone and nonmonotone circuit complexity is
  exponential.
\newblock {\em Combinatorica}, 8(1):141--142, 1988.

\bibitem[Urq87]{DBLP:journals/jacm/Urquhart87}
Alasdair Urquhart.
\newblock Hard examples for resolution.
\newblock {\em J. {ACM}}, 34(1):209--219, 1987.

\bibitem[Val80]{DBLP:journals/tcs/Valiant80}
Leslie~G. Valiant.
\newblock Negation can be exponentially powerful.
\newblock {\em Theor. Comput. Sci.}, 12:303--314, 1980.

\bibitem[Zhu17]{DBLP:conf/focs/Zhuk17}
Dmitriy Zhuk.
\newblock A proof of {CSP} dichotomy conjecture.
\newblock In Chris Umans, editor, {\em 58th {IEEE} Annual Symposium on
  Foundations of Computer Science, {FOCS} 2017, Berkeley, CA, USA, October
  15-17, 2017}, pages 331--342. {IEEE} Computer Society, 2017.

\end{thebibliography}
\end{document}
